\chapter{Introdução}
\label{chap:introducao}

\section{Motivação e contexto observacional}

As ondas gravitacionais permitem investigar a gravitação por meio de
sua propagação e de seus efeitos sobre sistemas de medição. Na
relatividade geral, as perturbações radiativas do vácuo possuem duas
polarizações tensoriais. Teorias alternativas podem modificar a relação
de dispersão, introduzir setores adicionais ou alterar simultaneamente
esses dois aspectos. A interpretação de um teste observacional exige,
portanto, relacionar a teoria adotada à resposta do detector e ao
procedimento estatístico que transforma os dados em limites ou
evidências entre modelos. Uma assinatura permitida pelas equações não é,
por si só, uma previsão de amplitude observável.

As redes de temporização de pulsares, designadas pela sigla PTA, do
inglês \emph{pulsar timing array}, investigam ondas de frequências muito
baixas por meio de resíduos nos tempos de chegada dos pulsos. Uma
perturbação gravitacional comum produz uma estrutura de correlação entre
pulsares que depende de sua distribuição angular, das polarizações e da
propagação. Para ondas dispersivas, essa estrutura também pode variar
com a frequência. A resposta incorpora contribuições avaliadas na Terra
e nos pulsares, com fases associadas às distâncias e à direção de
propagação. Esses elementos tornam as PTAs especialmente adequadas ao
estudo conjunto de massa e polarização, mas impedem identificar a
resposta observacional apenas com a deformação local de partículas
livres \autocite{liang2021,cordes2025}.

O interesse pelas baixas frequências decorre diretamente da relação de
dispersão de uma onda livre massiva no regime linear. Denotando por
\(m_g\) a massa do gráviton, por \(k\) o módulo do vetor de onda e por
\(h_{\mathrm P}\) a constante de Planck, adota-se
\begin{equation}
  \omega^2=c^2k^2+\left(\frac{m_gc^2}{\hbar}\right)^2,
  \qquad f_g=\frac{m_gc^2}{h_{\mathrm P}},
  \qquad \beta(f)=\frac{ck}{\omega}
  =\sqrt{1-\left(\frac{f_g}{f}\right)^2}.
  \label{eq:intro-dispersao}
\end{equation}
A velocidade de grupo é \(c\beta\), enquanto a velocidade de fase é
\(c/\beta\). Para uma massa fixa, o desvio cinemático cresce quando a
frequência se aproxima de \(f_g\). Abaixo desse limiar, a solução livre
considerada não descreve uma onda viajante com número de onda real.
Logo, a faixa de PTA oferece um regime qualitativamente diferente do
limite quase nulo, sem que isso garanta sensibilidade superior em toda
análise: espectro, ruído e tempo de observação também determinam a
informação disponível \autocite{liang2021,wu2024}.

Estudos recentes com dados do NANOGrav examinam tanto a dispersão massiva
quanto polarizações transversais adicionais
\autocite{wu2024,nanograv2024}. A interpretação dos resultados depende do
modelo de sinal, do suporte espectral e das hipóteses de ruído. Em
particular, a busca de polarização escalar transversal não apresentou
evidência significativa para sua inclusão \autocite{nanograv2024}.
Testes de dispersão de sinais de binárias realizados pelas colaborações
LIGO, Virgo e KAGRA oferecem outro contexto observacional, com hipóteses
de fonte e frequências distintas \autocite{lvk2026}. A presente
dissertação concentra-se na validade da inferência em PTAs, sem tratar
limites obtidos por métodos diferentes como medidas diretamente
intercambiáveis.

\section{Continuidade com o trabalho de graduação}

O ponto de partida é o trabalho de graduação de Wayne Leonardo Silva de
Paula, apresentado ao Instituto Tecnológico de Aeronáutica em 2003
\autocite{tg2003}. O estudo investigou os estados de polarização de ondas
gravitacionais com gráviton massivo na formulação de Visser,
empregando equações linearizadas, tetradas e componentes de curvatura.
A comparação entre frequências destacou a relevância potencial do
regime de baixas frequências. Seu núcleo científico foi publicado no
ano seguinte por de Paula, Miranda e Marinho
\autocite{depaula2004}. Essa publicação delimita a continuidade e o ponto
de partida fundamental desta dissertação: resgatar a formulação pioneira
desenvolvida no ITA e expandi-la com rigor analítico e computacional,
atualizando-a para o cenário observacional contemporâneo das redes de
temporização de pulsares.

Uma atualização teórica requer distinguir seis deformações geométricas
possíveis de seis graus de liberdade físicos independentes. A matriz
simétrica de marés \(E_{ij}=R_{0i0j}\) possui seis componentes. Os
vínculos das equações de campo, entretanto, podem relacionar essas
componentes. No modelo linear de Fierz--Pauli em quatro dimensões, um
campo massivo de spin dois possui cinco graus de liberdade; a
helicidade zero pode produzir deformações escalar transversal e
longitudinal vinculadas. Atribuir duas amplitudes estatisticamente
independentes a essas deformações mudaria o modelo físico
\autocite{liang2021}.

A atualização também exige delimitar as aproximações. O uso de uma
tetrada nula não torna nula a propagação da onda. Relações obtidas sob
hipóteses quase nulas devem ter seu domínio verificado antes de serem
aplicadas perto do limiar massivo. Expressões exatas e discussões dessas
limitações já estão disponíveis na literatura \autocite{hyun2019}.
Além disso, a análise de consistência física exige fundamentar o setor
massivo em uma teoria de campo bem estabelecida. Adota-se aqui a formulação
relativística de Fierz--Pauli como a referência linear consistente para o
gráviton massivo de spin dois, examinando as extensões não lineares de dRGT
e Hassan--Rosen como fundamentação teórica de fundo
\autocite{visser1998,park2010,derham2011,hassan2012}. Essa escolha assegura
que as propriedades de polarização e os vínculos entre modos físicos sejam
rigorosamente respeitados na resposta do detector.

Um princípio metodológico central orienta a interpretação dos limites físicos
obtidos: demonstrar analiticamente que um limite superior pode emergir como mero
artefato da priori na ausência de informação nos dados. Sob priori uniforme na
massa adimensional $u=f_gT\in[0,1]$ e verossimilhança não-informativa constante,
o quantil posterior de 95\% é identicamente $0,95$, gerando formalmente
$m_gc^2\leq0,95h_{\mathrm P}/T$ com ganho de informação nulo ($D_{\rm KL}=0$).
Para superar essa armadilha interpretativa, esta dissertação estabelece um
critério conjunto rigoroso: avaliar simultaneamente limites superiores de massa,
calibração da verossimilhança e ganho informacional efetivo frente à priori.

\section{Problema de pesquisa}

Parte dos produtos de PTA é apresentada na forma de correlações
angulares já comprimidas em frequência. Essa representação é útil para
visualizar a correlação espacial, mas sua interpretação por um modelo
dispersivo exige conhecer como a compressão foi construída. A previsão
para um estimador comprimido deve conservar os pesos espectrais, a
janela de observação e as hipóteses que participaram de sua definição.
Substituir a resposta dependente da frequência por uma avaliação em uma
única frequência introduz uma aproximação adicional à compressão.

O trabalho de Zhao e Wang sobre MeerKAT e previsões para SKA constitui
uma motivação direta: diante de produtos comprimidos, adota uma
frequência de referência associada ao inverso da duração observacional
e explicita limitações relativas à covariância disponível
\autocite{zhao2026}. Outros estudos recentes investigam correlações
massivas, múltiplas polarizações, distâncias finitas e variabilidade
entre realizações do fundo \autocite{choi2025,cordes2025,wu2025}.
Diante desse panorama dinâmico, o desafio científico central reside em
estabelecer um controle sistemático sobre a cadeia de inferência estatística:
quantificar, com rigor probabilístico e metodológico, as consequências
reais das simplificações espectrais e das aproximações de verossimilhança
amplamente empregadas na literatura sobre os limites físicos derivados.

A pergunta central é: \emph{em quais condições a avaliação da resposta
em uma frequência de referência preserva a inferência da massa do
gráviton, quando comparada à análise com frequência explícita e à
compressão aplicada de modo consistente ao modelo?} A extensão escalar
investiga se as mesmas escolhas alteram a recuperação de uma
contribuição de helicidade zero fisicamente vinculada.

Para precisar a comparação, considere esquematicamente o valor esperado
de um estimador de correlação entre os pulsares \(a\) e \(b\):
\begin{equation}
  \mu_{ab}(m_g)=\int df\,W_{ab}(f)S_h(f)
  \Gamma_{ab}(f;m_g).
  \label{eq:intro-compressao}
\end{equation}
Nessa expressão, \(S_h\) descreve o espectro do sinal,
\(\Gamma_{ab}\) representa a resposta correlacionada e \(W_{ab}\)
reúne os fatores do estimador e da observação, inclusive conversões de
unidades pertinentes. Uma aproximação de frequência única substitui
\(\Gamma_{ab}(f;m_g)\) por
\(\Gamma_{ab}(f_{\mathrm{ref}};m_g)\) na integral. A igualdade entre
as duas previsões precisa ser demonstrada no regime analisado; não
decorre apenas da escolha \(f_{\mathrm{ref}}=1/T\).

A hipótese de trabalho é que a adequação dessa aproximação depende da
variação da resposta ao longo da banda efetivamente medida e da
relevância dessa variação diante das incertezas. Uma banda estreita ou
uma resposta quase constante pode tornar as análises equivalentes na
precisão disponível. Uma resposta variável pode produzir discrepâncias
mensuráveis, mas a existência, o sinal e a magnitude de eventual viés
não são pressupostos. Também é possível perder informação com uma
compressão consistente sem introduzir o mesmo erro de modelagem da
aproximação de frequência única.

Essa separação exige ainda controlar a forma probabilística da verossimilhança.
Embora a literatura frequentemente assuma que a calibração de médias e covariâncias
seja suficiente para validar uma aproximação gaussiana, formas quadráticas de dados
complexos possuem assimetria e caudas pesadas que podem comprometer a cobertura
estatística. Portanto, a investigação proposta desacopla metodologicamente três
níveis de aproximação: a hipótese de normalidade da verossimilhança, a contração
linear do espaço de dados (compressão) e a substituição da resposta dependente da
frequência por um template congelado.

Esse desenho experimental preenche uma \textbf{lacuna metodológica essencial} na
literatura de ondas gravitacionais. Embora estudos antecedentes tenham abordado
previsões multifrequência \autocite{han2026} ou discutido a não-gaussianidade de
fundos estocásticos \autocite{franciolini2025}, as análises observacionais aplicadas
à gravidade dispersiva rotineiramente combinam essas três operações sem isolar a
origem de eventuais vieses ou perdas de sensibilidade. A contribuição original
desta dissertação consiste na concepção de um arcabouço diagnóstico pareado com
um controle gaussiano exato ($G$), permitindo quantificar matematicamente a perda
líquida de informação, auditar a calibração de verossimilhanças por SBC com
tratamento rigoroso de incertezas numéricas e mapear os limites operacionais de
validade das aproximações em frequência.

\section{Objetivos e critérios de avaliação}

O objetivo geral é construir uma comparação reproduzível entre três
formas de inferência em PTAs: frequência explícita, compressão
consistente e resposta avaliada em uma frequência de referência. As
três análises partem das mesmas realizações simuladas e de
hipóteses físicas compatíveis, permitindo atribuir diferenças ao
tratamento da frequência e às aproximações declaradas.

Os objetivos específicos são:
\begin{enumerate}
  \item Reproduzir o núcleo linear do TG, documentando convenções,
  unidades, vínculos e limites de validade, com verificações algébricas
  independentes dos resultados estatísticos posteriores.
  \item Implementar a resposta tensorial dispersiva de PTA com termos
  da Terra e dos pulsares, validando integrações e limites analíticos.
  \item Gerar conjuntos sintéticos com massa, espectro, geometria,
  distâncias e ruído conhecidos, preservando a covariância necessária à
  comparação dos estimadores.
  \item Quantificar diferenças entre as três análises por meio de
  recuperação de parâmetros, incertezas, calibração e informação
  fornecida pelos dados em relação à priori.
  \item Avaliar a robustez às escolhas a priori, ao suporte espectral e
  ao tratamento das covariâncias, separando restrições cinemáticas de
  informação observacional.
  \item Estender cenários selecionados à helicidade zero de
  Fierz--Pauli, mantendo o vínculo entre suas deformações transversal e
  longitudinal, e preparar resultados reproduzíveis para um artigo.
\end{enumerate}

A avaliação não ficou restrita à concordância visual de curvas. A
verificação física buscou recuperar o limite tensorial sem massa nas
condições pertinentes; a verificação numérica avaliou
convergência em precisão suficiente para a inferência. Campanhas de
injeção e recuperação mediram deslocamentos de parâmetros e
incertezas. Calibração sob a priori e cobertura em parâmetros fixos
foram avaliadas separadamente, pois respondem a perguntas distintas.
Quando houver comparação entre modelos, a distribuição da estatística
sob a hipótese nula foi medida antes de interpretar uma preferência
como evidência de física adicional.

\section{Delimitação e viabilidade}

O núcleo do trabalho compreende
propagação linear sobre fundo plano, resposta de PTA, simulações e
inferência estatística. As amplitudes dos setores considerados serão
hipóteses de população explicitadas. Não serão deduzidas a emissão
completa de binárias em uma teoria não linear, a história cosmológica
do fundo ou uma solução do mecanismo de Vainshtein. Essa delimitação
permite testar a inferência sem apresentar um modelo fenomenológico de
amplitudes como previsão universal de gravidade massiva.

A estratégia de investigação baseia-se em simulações sintéticas de alta fidelidade e benchmarks numéricos rigorosamente controlados, incorporando as coordenadas celestes reais de pulsares do catálogo público MeerKAT. Esse desenho deliberado garante total controle sobre a verdade injetada, permitindo isolar causalmente os efeitos estatísticos e cinemáticos sem as ambiguidades de modelos empíricos de ruído ou amostragem temporal irregular. Paralelamente, o trabalho investiga os fundamentos da transição para dados reais de temporização (TOAs), demonstrando formalmente no Capítulo~\ref{cap:c11-parcial} como a matriz de ajuste de timing induz correlações entre frequências e pseudocovariância, estabelecendo assim os pré-requisitos exatos para a transição observacional.

\section{Organização da dissertação}

O Capítulo~\ref{chap:revisao} apresenta a revisão bibliográfica e delimita a contribuição em relação aos antecedentes. O Capítulo~\ref{chap:fundamentos} reúne os fundamentos, a reconstrução do trabalho original e os vínculos de Fierz--Pauli. A resposta tensorial de PTA é derivada e validada no Capítulo~\ref{chap:resposta-pta}, e o experimento sintético é definido no Capítulo~\ref{chap:simulacao}.

O Capítulo~\ref{chap:validacao} descreve a inferência e seus controles. A comparação de compressões é apresentada no Capítulo~\ref{chap:compressao}, seguida da análise de robustez no Capítulo~\ref{chap:robustez}. A extensão de helicidade zero e a população ampliada complementam os testes. A discussão e as conclusões integram os resultados, distinguindo limitações físicas, estatísticas e numéricas.
